\section{Related Work}
\label{sec:related}

\subsection{Performance- and network-aware microservice scheduling}

A substantial body of research has addressed microservice placement and
orchestration by optimizing communication cost, response time, throughput,
or resource utilization. Rodrigues et al. formulated network-aware
container scheduling as a multi-criteria optimization problem that jointly
considers container requirements and network-level constraints
\citep{rodrigues2020networkaware}. Their formulation demonstrates that
placement decisions should account for communication requirements in
addition to compute availability. However, it targets network-efficient
container placement rather than online admission of sensitive services
under policy and near-term operational-risk constraints. Wojciechowski et al. proposed NetMARKS, a Kubernetes scheduler that uses
service-mesh telemetry to estimate communication between a pod awaiting
placement and the pods already assigned to each node
\citep{wojciechowski2021netmarks}. By favouring destinations that reduce
cross-node service traffic, NetMARKS directly incorporates application
communication into pod placement. Its scoring mechanism is nevertheless
centred on traffic locality and does not distinguish policy-inadmissible
nodes, predict next-interval node overload, or regulate repeated
destination use. Shen et al. developed KaiS, a collaborative learning-based scheduling
framework for Kubernetes-oriented edge--cloud networks
\citep{shen2023kais}. KaiS combines coordinated multi-agent actor--critic
learning for request dispatch with graph representation learning and
multiple policy networks for service orchestration. Its two-timescale
design jointly adapts request routing and service placement to improve
long-term request-processing throughput. KaiS is an important example of
learning-based orchestration over heterogeneous edge and cloud resources,
but its learned policy optimizes performance over the available
infrastructure rather than treating jurisdiction, trust, encryption, or
hardware-isolation requirements as hard placement gates. Song et al. presented ChainsFormer for chain-latency-aware microservice
resource provisioning \citep{song2023chainsformer}. The method identifies
critical service chains and microservices and applies reinforcement
learning to coordinate vertical and horizontal scaling. ChainsFormer
therefore captures dependencies that are invisible to isolated
microservice scaling. Its action space is resource provisioning and
replica adjustment, whereas CAGE-S decides whether a current placement
should be retained and which policy-admissible destination should be
selected. Xiangchun Chen et al. proposed DTDRLMO, which combines digital twins and
deep reinforcement learning for joint microservice offloading and
bandwidth allocation in collaborative edge computing
\citep{chen2023dtdrlmo}. Digital-twin models are used to anticipate changes
in edge-node loads and network conditions, after which the learned policy
selects offloading destinations to minimize average service-completion
time. DTDRLMO is particularly relevant because it introduces predicted
environment state into placement. It does not, however, formulate an
explicit next-interval threshold-crossing probability, enforce
service-specific policy admissibility, or apply destination-use fairness
and a calibrated pre-commit headroom check. Ming Chen et al. introduced TraDE, a network- and traffic-aware adaptive
scheduler that redeploys microservices when changing traffic and
inter-node delays degrade application QoS \citep{chen2026trade}. TraDE
constructs a bidirectional service-traffic representation and a node-delay
matrix, and uses a parallel greedy mapper to determine a new deployment.
It is one of the closest recent works in terms of dynamic microservice
redeployment. Its scheduling signal is nevertheless derived mainly from
service traffic, communication delay, load, and application performance.
CAGE-S addresses a complementary decision problem in which service
legality first defines the candidate space, and temporal node risk,
historical destination use, and projected headroom determine the final
placement within that space. Taken together, these methods provide sophisticated mechanisms for
traffic-aware placement, adaptive offloading, resource provisioning, and
learning-based orchestration. Their principal objective is to improve
application performance or infrastructure efficiency. They do not
explicitly organize hard policy admission, next-interval node-risk
prediction, destination-use fairness, and commit-time resource
verification as consecutive and non-interchangeable stages.

\subsection{Telemetry-driven SLO and resource management}

A second line of work uses runtime telemetry to identify performance
degradation and adjust microservice resources. Qiu et al. developed FIRM,
which combines fine-grained telemetry with machine learning to detect
service-level objective violations, localize contended resources, and
reprovision affected microservices \citep{qiu2020firm}. FIRM shows that
microservice resource management benefits from interpreting application
and infrastructure signals jointly. Its objective is corrective resource
reallocation after identifying performance interference rather than
policy-constrained destination selection. Gan et al. proposed Sage for diagnosing unpredictable performance in
interactive cloud microservices \citep{gan2021sage}. Sage uses
unsupervised learning and service-dependency information to identify the
root causes of QoS violations and apply corrective actions. The method
addresses the difficult problem of locating performance faults within a
microservice graph. It does not define which nodes are legally admissible
for a sensitive service or control the fairness of repeated placement
decisions. Zhang et al. introduced Sinan, a data-driven QoS-aware cluster manager
that models dependencies among microservice tiers and allocates resources
to satisfy tail-latency targets \citep{zhang2021sinan}. Sinan demonstrates
that instantaneous utilization alone is insufficient for managing
end-to-end microservice performance and that historical traces can improve
resource decisions. Its decisions concern per-tier resource allocation,
not legality-first placement across nodes with heterogeneous governance
and trust attributes. Fettes et al. developed Reclaimer, a deep-reinforcement-learning approach
for dynamically assigning CPU cores to cloud microservices
\citep{fettes2023reclaimer}. Reclaimer adapts to changes in the number and
behaviour of microservices while seeking to reduce resource allocation
without violating application QoS. This resource-allocation objective is
different from CAGE-S, where the primary action is retaining or changing a
service destination under hard admission conditions. Wen et al. proposed StatuScale, which detects workload status and combines
vertical and horizontal scaling to respond to request spikes and bursts
\citep{wen2024statuscale}. The status detector determines which scaling
strategy should be activated, while a horizontal controller manages the
number of service replicas. StatuScale is relevant to CAGE-S because both
methods distinguish evolving resource behaviour from a single static
snapshot. StatuScale uses this information to choose a scaling action;
CAGE-S uses temporal node history to estimate next-interval destination
risk before placement commitment. Hu et al. introduced LSRAM, a lightweight autoscaling and SLO resource
allocation framework based on gradient descent
\citep{hu2024lsram}. LSRAM rapidly allocates service-level resource
budgets and updates its allocation model when workloads, applications, or
cluster conditions change. Its lightweight adaptation reduces the cost of
retraining complex resource models. However, it addresses resource
allocation across a service chain rather than policy-aware node
admissibility and destination ranking. These systems establish the importance of online telemetry, dependency
structure, temporal context, and adaptive control. Their action spaces
mainly involve diagnosis, reprovisioning, CPU allocation, or replica
scaling. CAGE-S instead uses temporal information inside a node-placement
pipeline: history estimates the probability of a destination crossing an
operational boundary, while a separate Q95 check verifies projected CPU
and memory headroom before commitment. This separation prevents the
predictive model from serving simultaneously as both a risk estimator and
an executability rule.

\subsection{Fairness, policy, and trusted edge orchestration}

Fairness has also been studied in edge scheduling, although usually at a
different level from destination fairness. Madej et al. proposed
client-fair, priority-fair, and hybrid scheduling policies for
resource-constrained edge nodes \citep{madej2020fair}. Their work balances
the allocation of edge execution opportunities among clients and job
priorities. CAGE-S addresses a complementary form of fairness: it uses
historical destination-selection counts to prevent repeated migrations
from concentrating on a small set of nodes that consistently receive
favourable operational scores. Pallewatta et al. developed a comprehensive taxonomy of microservice-based
IoT application placement in edge and fog environments
\citep{pallewatta2023taxonomy}. Their analysis identifies heterogeneity,
application composition, placement dynamics, security, privacy, and
evaluation methodology as central dimensions of the problem. The taxonomy
also highlights that microservice placement is not determined by resource
capacity alone, particularly when applications span multiple
administrative and infrastructure domains. Pallewatta and Babar subsequently proposed a Policy-as-Code architecture
for secure management of edge--cloud IoT microservices
\citep{pallewatta2024policy}. Their framework uses Kubernetes, Istio, and
Open Policy Agent to enforce security policies during initial placement,
scaling, migration, and dynamic service composition. This work is the
closest to CAGE-S in its treatment of security policy as an explicit
control-plane concern. Its primary contribution is policy specification
and enforcement; it does not rank policy-admissible nodes using calibrated
next-interval operational risk, destination-use history, and an empirical
resource guardrail. Thijsman et al. presented a remote-attestation architecture for enrolling
and authorizing Kubernetes workers at the edge
\citep{thijsman2024attested}. Their design uses hardware-rooted boot
attestation and dynamically grants or revokes credentials and
role-based-access-control privileges according to attestation events.
This mechanism provides a strong foundation for determining whether an
edge worker can be trusted. It does not itself decide which trusted worker
is operationally preferable for a particular service. In CAGE-S,
attestation or trusted-execution status can contribute to the policy mask,
after which temporal risk, fairness, and pre-commit headroom determine the
destination. The policy and trust literature therefore provides mechanisms for
expressing mandatory requirements and establishing trustworthy cluster
membership, while the fairness literature prevents inequitable access or
allocation. CAGE-S connects these concerns at the scheduling-decision
level: policy and trust determine admissibility, historical destination
use controls concentration, and telemetry-derived risk determines
operational suitability.

\subsection{Positioning of CAGE-S}

Table~\ref{tab:related-positioning} compares CAGE-S with the representative
systems reviewed above. The comparison records only capabilities that are
explicit components of each published method. General resource-capacity
constraints, load penalties, or predicted environment states are not
classified as a calibrated pre-commit guardrail or an explicit
next-interval threshold-risk model.

\begin{table*}[t]
\centering
\caption{Positioning of CAGE-S against representative related systems.
``Explicit risk'' denotes prediction of next-interval node-threshold
crossing, and ``pre-commit check'' denotes a separate resource-headroom
verification immediately before placement commitment.}
\label{tab:related-positioning}
\scriptsize
\setlength{\tabcolsep}{3pt}
\begin{adjustbox}{max width=\textwidth}
\begin{tabular}{
p{2.65cm}
p{2.75cm}
p{2.15cm}
ccccc
p{3.15cm}}
\toprule
System &
Primary action &
Dynamic information &
Policy/
trust gate &
Explicit
risk &
Fair-use
control &
Pre-commit
check &
Online
replacement &
Primary objective \\
\midrule

Rodrigues et al.
\citep{rodrigues2020networkaware} &
Container placement &
Network and container requirements &
No &
No &
No &
No &
No &
Multi-criteria network-aware placement \\

NetMARKS
\citep{wojciechowski2021netmarks} &
Kubernetes pod placement &
Service-mesh traffic &
No &
No &
No &
No &
No &
Inter-service traffic reduction \\

KaiS
\citep{shen2023kais} &
Request dispatch and orchestration &
Online system and graph state &
No &
No &
No &
No &
Yes &
Long-term request throughput \\

ChainsFormer
\citep{song2023chainsformer} &
Vertical and horizontal scaling &
Chain latency and workload state &
No &
No &
No &
No &
No &
Latency and resource efficiency \\

DTDRLMO
\citep{chen2023dtdrlmo} &
Offloading and bandwidth allocation &
Predicted node and network state &
No &
No &
No &
No &
Yes &
Average service-completion time \\

TraDE
\citep{chen2026trade} &
Adaptive microservice redeployment &
Traffic stress and node delay &
No &
No &
Partial &
No &
Yes &
Response time and goodput \\

FIRM
\citep{qiu2020firm} &
Resource reprovisioning &
Resource and application telemetry &
No &
No &
No &
No &
No &
SLO recovery and utilization \\

Sage
\citep{gan2021sage} &
Diagnosis and corrective action &
Traces and service dependencies &
No &
No &
No &
No &
No &
Root-cause detection and QoS recovery \\

Sinan
\citep{zhang2021sinan} &
Per-tier resource allocation &
Online traces and dependency state &
No &
No &
No &
No &
No &
Tail latency and resource efficiency \\

Reclaimer
\citep{fettes2023reclaimer} &
Dynamic CPU allocation &
Runtime microservice state &
No &
No &
No &
No &
No &
QoS with reduced CPU allocation \\

StatuScale
\citep{wen2024statuscale} &
Vertical and horizontal scaling &
Detected workload status &
No &
No &
No &
No &
No &
Burst-aware elastic scaling \\

LSRAM
\citep{hu2024lsram} &
Autoscaling and SLO allocation &
Adaptive cluster and workload state &
No &
No &
No &
No &
No &
QoS and resource efficiency \\

Madej et al.
\citep{madej2020fair} &
Edge-job scheduling &
Queue and job priorities &
No &
No &
Yes &
No &
No &
Client and priority fairness \\

Policy-as-Code
\citep{pallewatta2024policy} &
Secure placement, scaling, and migration &
Service and infrastructure policies &
Yes &
No &
No &
No &
Yes &
Security-policy enforcement \\

Attested workers
\citep{thijsman2024attested} &
Worker enrollment and authorization &
Hardware-attestation events &
Yes &
No &
No &
No &
No &
Trusted cluster membership \\

\cage &
Source retention or migration &
Current and historical node telemetry &
Yes &
Yes &
Yes &
Yes &
Yes &
Compliance, temporal risk, fairness,
and projected headroom \\
\bottomrule
\end{tabular}
\end{adjustbox}
\end{table*}

Among the representative systems in
Table~\ref{tab:related-positioning}, CAGE-S is the only scheduler that
explicitly combines all four decision responsibilities considered in this
work. First, a hard policy gate determines whether a node may host the
service. Second, a calibrated temporal model estimates whether an
admissible node is likely to cross the operational boundary during the
next control interval. Third, historical destination use limits repeated
placement concentration. Fourth, an independent Q95 guardrail verifies
projected resource headroom before commitment. The distinction is therefore not merely the use of a different predictive
model. CAGE-S changes the structure of the scheduling decision by assigning
permission, near-term operational risk, distribution of repeated choices,
and commit-time executability to separate ordered stages. This design
prevents favourable performance scores from compensating for mandatory
policy failures and prevents a predictive probability from replacing a
direct pre-commit resource check.